\section{Constitutive, Geometric, Boundary, and Observable Modeling}
\label{sec:modeling-closures}
\label{sec:modeling-closures-observables}

Choosing a model tier does not close a hemodynamic problem. The mathematical
state must be paired with constitutive, geometric, wall, boundary, and
observable conventions. These choices determine which terms appear in the
equations and which reported quantities are native states, derived fields,
diagnostics, or cross-tier observation maps.

\subsection{Constitutive and Wall Closures}
\label{subsec:constitutive-wall-closures}

Blood is often modeled as an incompressible Newtonian fluid in large-vessel
simulations, especially when the target quantities are dominated by moderate to
high shear rates and the uncertainty in geometry or boundary data is already
large. Non-Newtonian categories remain important in stenotic, small-vessel, or
low-shear regimes because viscosity may depend on shear rate, hematocrit, and
diameter \cite{PriesEtAl1992BloodViscosityTubeFlow}. In a review taxonomy,
Carreau and Carreau--Yasuda laws represent shear-thinning models with limiting
viscosities, Casson laws introduce a yield-stress-like response, and power-law
models represent scale-free shear-rate dependence. Computational studies of
Carreau--Yasuda and intracranial stenosis rheology illustrate why the selected
rheology must be reported with the resulting hemodynamic observable
\cite{KumarEtAl2019CarreauYasudaHemodynamics,LiuEtAl2021ICASRheology}.

Wall modeling is a separate closure axis. A fixed-wall model prescribes the
lumen boundary and solves only the fluid problem. A reduced compliant-wall model
uses an algebraic or differential pressure--area relation to close a 1D or
network state. A resolved FSI model solves structural wall motion and couples it
to the fluid at the interface
\cite{FormaggiaEtAl2003OneDimensionalBloodFlow,FormaggiaEtAl2007FSICoupling}.
The same rheology can be paired with any of these wall choices, so rheology and
wall mechanics should not be merged into a single modeling label.

\subsection{Geometry, Boundary Data, and Observables}
\label{subsec:geometry-boundary-observable-closures}

A stenosis can enter the mathematics as a resolved lumen narrowing, an
axisymmetric radius profile, a 1D reference-area variation, an empirical
pressure-loss relation, or a network resistance. Distributed 1D formulations
make the reference area and its derivatives part of the balance-law closure,
while resolved models represent the same narrowing through the computational
domain. Extended 1D stenosis formulations show that variable-radius effects may
require more than simply replacing a healthy radius with a narrowed radius in a
standard area--flow model \cite{CanicEtAl2024Extended1DStenosis}.

Boundary data are equally model dependent. Inlets may be prescribed as
velocity profiles, mean-flow waveforms, pressure histories, or characteristic
incoming data. Outlets may be prescribed as pressure, traction, resistance,
Windkessel, impedance, reflection, or characteristic conditions. In multiscale
models, the same words also name interface contracts between tiers, not only
external boundary conditions
\cite{QuarteroniVeneziani2003GeometricalMultiscale,ShiEtAl2011BloodFlowModelReview}.

Observables should be named at the level where they are actually defined.
Pressure may be an absolute pressure, gauge pressure, cross-sectional mean
pressure, diagnostic pressure reconstructed from a wall law, or terminal
compartment pressure. Flow may be a surface integral, a 1D state variable, or a
network edge variable. Velocity may be a resolved vector field, a section
average, an assumed radial profile, or a sampled centerline quantity. WSS is a
traction-derived wall quantity requiring stress and wall geometry, not a generic
name for any reduced shear proxy \cite{JohnEtAl2017WallShearStressVectorForm}.
Displacement belongs to a wall or structural state, so it is native to resolved
FSI or a structural model and only indirectly represented by a reduced
pressure--area wall law.

\begin{table}[!htb]
    \centering
    \scriptsize
    \caption{Closure and observable axes that should remain explicit when
    comparing stenotic-flow models.}
    \label{tab:modeling-closures-observables}
    \begin{tabularx}{\textwidth}{@{}
        >{\raggedright\arraybackslash}p{0.20\textwidth}
        >{\raggedright\arraybackslash}p{0.34\textwidth}
        >{\raggedright\arraybackslash}X@{}}
        \toprule
        \textbf{Axis}
            & \textbf{Review categories}
            & \textbf{Observable consequence} \\
        \midrule
        Rheology
            & Newtonian; Carreau; Carreau--Yasuda; Casson; power-law; tube-flow
            effective-viscosity laws.
            & Changes viscosity, friction, shear-rate interpretation, and sometimes
            pressure or WSS trends. \\
        Wall mechanics
            & Fixed wall; algebraic pressure--area law; compliant 1D wall;
            resolved FSI wall.
            & Determines whether area, pressure, displacement, and interface work are
            solved states or reconstructed quantities. \\
        Stenosis geometry
            & Resolved surface; axisymmetric radius; 1D reference-area profile;
            empirical loss; terminal/network resistance.
            & Determines whether throat effects appear through geometry, source
            terms, pressure loss, or lumped parameters. \\
        Inlet data
            & Resolved velocity profile; mean-flow waveform; pressure waveform;
            characteristic incoming data.
            & Fixes what enters the domain and whether section averages or local
            profiles are prescribed. \\
        Outlet data
            & Pressure, traction, resistance, Windkessel, impedance, reflection, or
            characteristic condition.
            & Controls wave reflection, terminal load, and pressure--flow phase. \\
        Observation map
            & Pressure, flow, velocity, WSS, displacement, or derived diagnostic.
            & Separates native model states from reconstructed or cross-tier
            comparison quantities. \\
        \bottomrule
    \end{tabularx}
\end{table}

\begin{stenosisillustration}
\label{ex:stenosisillustration-modeling-closures}
In the StenosisHemodynamics case study, the principal reported problem is a
single-vessel distributed 1D stenosis model. The active review placement is:
Newtonian viscosity for the principal comparison; a parabolic velocity profile
for the main area--flow closure; an elastic pressure--area wall law; a smooth
analytic stenosis reference radius; a prescribed reference-area inlet flow; and
a fixed-area characteristic outlet approximation. The native states are reduced
area and flow, with mean velocity and pressure obtained through the declared
solver-coordinate and wall-law maps. The resolved-data comparison therefore
reads as a velocity-observation comparison, not as a direct WSS or
wall-displacement comparison.
\end{stenosisillustration}
